% ===================== OBJETIVOS =====================
\section{Objetivos del Proyecto}
\label{sec:objetivos}

\subsection{Objetivo General}

Evaluar la viabilidad hidráulica y los riesgos operativos de la reubicación de dos tuberías de \SI{12}{in} NPS en HDPE PE100 SDR 17 (salmuera saturada y condensado) mediante Perforación Horizontal Dirigida (PHD) en el Humedal Arrieros, comparando el impacto dinámico y estático con respecto al trazado superficial original para garantizar la operación segura y la presión de succión requerida en la estación Km8 de BRINSA S.A.

\subsection{Objetivos Específicos}

Los objetivos específicos del presente estudio se orientan a cuantificar las variables hidráulicas y mecánicas del sistema, descomponiéndose en las actividades listadas en la Tabla~\ref{tab:obj_especificos}.

\begin{table}[htbp]
    \centering
    \caption{Objetivos específicos para la evaluación del diseño hidráulico.}
    \label{tab:obj_especificos}
    \setcounter{itemcount}{0}
    \begin{tabularx}{\textwidth}{@{}NX@{}}
        \toprule
        \multicolumn{1}{c}{\textbf{Item}} & \textbf{Descripción} \\
        \midrule
        & Cuantificar y comparar las pérdidas de presión dinámicas por fricción en flujo continuo para el trazado superficial original y la alternativa de profundización dirigida. \\
        & Evaluar el impacto del cambio geométrico en la presión de succión disponible en la estación de bombeo de mantenimiento Km8. \\
        & Determinar la carga de presión hidrostática máxima en el fondo de la perforación (cota \SI{2542.75}{m.s.n.m.}) para validar la resistencia de la tubería PE100 SDR 17. \\
        & Identificar los riesgos de acumulación de sólidos y entrampamiento de gases en la geometría de sifón invertido de la PHD, proponiendo las acciones correctivas correspondientes. \\
        \bottomrule
    \end{tabularx}
\end{table}
